\subsection{Elicitazione dei Requisiti}
\label{elicitazione}
La seconda fase precedentemente descritta richiedere un maggiore approfondimento. Nella ricerca dei requisiti si può incappare in problemi relativi al linguaggio naturale o a processi mentali quali rimozione, distorsione, generalizzazione. \\
La scoperta dei requisiti avviene attraverso comunicazione diretta con gli attori, cercando di favorire la comunicazione con un gruppo eterogeneo piuttosto che omogeneo. A tale fine, si selezionando attori provenienti da gruppi divisi per:
\begin{itemize}
    \item \textbf{punto di vista}
    \begin{itemize}
        \item \textbf{primari}, interagiscono direttamente col sistema e attivano casi d'uso
        \item \textbf{finali}, l'uso del sistema da parte di terzi realizza i loro obiettivi
        \item \textbf{di supporto (o secondari)}, offrono un servizio al sistema
        \item \textbf{fuoriscena}, sono interessati al comportamento del sistema ma non ci interagiscono direttamente.
    \end{itemize}
    \item \textbf{ruolo}: persone, sistemi esterni, sensori
    \item \textbf{interfacce per l'interazione col sistema}
\end{itemize}
\subsection{Tecniche di elicitazione}
Per la scoperta dei requisiti si sfruttano specifiche tecniche. Le principali sono:
\paragraph{Focus group} Una specifica parte del team approfondisce tematiche del sistema poco chiari.
\paragraph{Osservazione} Si osservano i potenziali utilizzatori del sistema nello svolgimento delle regolari mansioni.
\paragraph{Questionari} Viene distribuita agli stakeholder una lista di domande per comprendere concetti specifici.
\paragraph{Perspective-based Reading} Identificazione e studio di documenti di rilievo per lo sviluppo.
\paragraph{Intervista} \hyperref[intervista]{Approfondimento.} 
\paragraph{Workshop} \hyperref[workshop]{Approfondimento.}
\subsection{Tecniche Accessorie}
\label{tecniche-accessorie}
Sono inoltre spesso usate, in combinazione con le tecniche sopra elencate, alcune tecniche accessorie quali:
\paragraph{Brainstorming} Emersione di idee innovative.
\paragraph{KJ Method} Brainstorming in contesti con partecipanti eterogenei.
\hyperref[kjmethod]{Approfondimento.}
\paragraph{Prototyping} Concretizzazione pragmatica di alcuni aspetti del sistema.
\paragraph{Mind Mapping} Costruzione di mappe concettuali.
\paragraph{Elicitation checklist} Lista di tipici aspetti in un particolare contesto applicativo.
\subsection{Intervista}\label{intervista}
L'\textbf{Intervista} è una tecnica in cui l'analista dei requisiti conversa con un gruppo di stakeholder per esplorare le sue necessità o aspettative e i vincoli a cui il sistema deve essere soggetto. \\
Ci sono tre tipi di interviste:
\begin{itemize}
    \item \textbf{standardizzate}, con domande chiuse
    \item \textbf{esplorative}, con domande aperte
    \item \textbf{destrutturate}, con una conversazione attiva guidata dall'analista
\end{itemize}
L'intervista si compone di più fasi:
\begin{itemize}
    \item \textbf{preparazione}, definizione delle modalità di intervista, delle domande da porre, delle possibili problematiche e degli obiettivi su cui chiedere chiarimenti, scelta degli intervistati ed individuazione del luogo e del tempo in cui svolgere l'intervista.
    \item \textbf{esecuzione}, l'effettiva esecuzione dell'intervista, si divide ulteriormente in:
    \begin{itemize}
        \item \textbf{apertura}, introduzione di obiettivi e motivazioni
        \item \textbf{conduzione}, l'analista pone le domande, prende appunti e chiede chiarimenti su risposte poco chiare
        \item \textbf{chiusura}, riassunto dei punti salienti emersi durante la conduzione
    \end{itemize}
    \item \textbf{rielaborazione}, analisi delle risposte raccolte per estrarre requisiti
\end{itemize}
La tecnica dell'intervista permette un'esplorazione profonda del punto di vista di uno stakeholder, comprese necessità nascoste o di cui c'è poca consapevolezza. E' però soggetta alla necessità di tempo per l'organizzazione, non garantisce l'individuazione di tutti i possibili requisiti ed è soggetto a bias, tra i quali:
\begin{itemize}
    \item \textbf{bias di selezione}, il gruppo selezionato dall'intervistatore potrebbe non avere un'idea completa degli obiettivi o dei limiti del sistema
    \item \textbf{bias di conferma}, l'intervistatore potrebbe inconsciamente porre domande che riconfermino idee o aspettative già formate
    \item \textbf{bias di interpretazione}, in fase di rielaborazione potrebbero sorgere ambiguità dovute ad errori di interpretazione
\end{itemize}
\subsection{Workshop}\label{workshop}
Il workshop è una tecnica di elicitazione basata sulla collaborazione tra un gruppo di stakeholder e l'analista. Si divide in tre fasi:
\begin{itemize}
    \item \textbf{apertura}, si descrivono obiettivi, regole e possibili tecniche ausiliarie usate
    \item \textbf{conduzione}, avviene similarmente a quella dell'intervista, ma nel workshop la figura dell'analista viene accompagnata a quella del minute-taker, che trascrive tutto ciò che viene affermato durante la riunione in appunti chiamati "minutes"
    \item \textbf{chiusura}, stabilita dall'analista, ne segue la raccolta dei risultati
\end{itemize}
In seguito alla chiusura, il \textbf{minute-taker} riordina gli appunti presi e richiede feedback dei vari partecipanti al workshop, modificando i risultati a seconda dei loro responsi. \\
Alla luce di ciò, la principale caratteristica del workshop è la partecipazione attiva di più stakeholder e il confronto diretto tra essi e il minute-taker. Inoltre, permette facilmente di affiancare varie \hyperref[tecniche-accessorie]{tecniche accessorie} e di raccogliere in modo efficace più prospettive, identificando anche eventuali conflitti tra le varie visioni degli stakeholder. D'altro canto, richiede un lavoro maggiore, in quanto è spesso richiesto di filtrare, ordinare ed analizzare idee emerse, oltre che richiedere una disponibilità costante da parte degli stakeholder.
\subsection{KJ Method}\label{kjmethod}
Il \textbf{KJ Method} è una tecnica accessoria usata per organizzare e raggruppare un gran numero di idee emerse. La tecnica consiste nella generazione di idee o input individuale (in genere trascritte su oggetti fisici come post-it o note) volte alla risoluzione di una problematica o al raggiungimento di un obiettivo. Una volta terminata la fase di riflessione individuale, le idee vengono esposte e lette dall'intero gruppo, si procede poi per raggruppamento delle idee simili, sintetizzandone con una breve etichetta il contenuto. Si votano i temi più rilevanti, così da darli priorità. \\
Il KJ Method permette l'emergere di requisiti lasciati implicitati, temi ricorrenti o ritenuti prioritari ed eventuali correlazioni tra esigenze diverse. Richiede, però, che siano disponibili input sufficienti da raggruppare e che sia presente una moderazione per evitare bias in fase di clustering. \\
Non rappresenta una tecnica di elicitazione completa, ma uno strumento di supporto con cui chiarire ed ordinare idee già emerse.